\section{Streaming ASR}

\begin{frame}[plain]
  \centering
  \vspace{0.10\textheight}
  {\LARGE \textbf{Streaming ASR}}

  \vspace{0.08\textheight}
  {\large Transcribing while the user is still speaking}
\end{frame}

\begin{frame}{Offline vs Streaming Recognition}
  \begin{itemize}
    \item \textbf{Offline (full-context):} the whole utterance is available before decoding starts. Best accuracy; fine for transcribing recordings.
    \item \textbf{Streaming:} tokens must appear while audio is still arriving. Required for voice assistants, live captions, and dictation.
    \item The core constraint: a streaming model may only use the \textbf{past} (and a small fixed lookahead), never the whole future of the utterance.
    \item User-perceived latency comes from three places: acoustic lookahead, model computation, and the decision that the speaker has finished.
  \end{itemize}
\end{frame}

\begin{frame}{Making the Encoder Streamable}
  \begin{itemize}
    \item Full self-attention encoders (Conformer, Whisper) attend over the entire utterance, so they cannot stream as-is.
    \item \textbf{Causal encoders:} restrict attention and convolutions to the past only. Cheapest latency, largest accuracy loss.
    \item \textbf{Chunked attention:} process audio in fixed chunks with a small right-context window; each chunk sees a little of the future.
    \item Chunk size is the latency-accuracy knob: bigger chunks recover offline accuracy, smaller chunks respond faster.
  \end{itemize}
\end{frame}

\begin{frame}{Why RNN-T Dominates Streaming}
  \begin{itemize}
    \item The RNN-Transducer is \textbf{frame-synchronous}: at each audio frame it either emits a token or a blank and moves on, so partial transcripts appear as audio arrives.
    \item No attention over the full input and no need to wait for the end of the utterance, unlike encoder-decoder models with cross-attention.
    \item The prediction network conditions on previously emitted tokens, giving language-model-like context without breaking streamability.
    \item Production systems (mobile dictation, on-device assistants) are overwhelmingly RNN-T variants with causal or chunked encoders.
  \end{itemize}
\end{frame}

\begin{frame}{Endpointing}
  \begin{itemize}
    \item \textbf{Endpointing} is deciding that the speaker has finished, so the final result can be committed and the microphone closed.
    \item Too eager: the system cuts users off mid-sentence. Too patient: every interaction feels slow.
    \item Signals used: silence duration, the decoder's end-of-utterance probability, and prosodic cues.
    \item Endpointing errors often dominate perceived quality even when word accuracy is high.
  \end{itemize}
\end{frame}
